\section{Protobuf Parser}\label{sec:proto-parse}

Following the two-step process outlined in Section~\ref{sec:philosophy}, the
Pollux Protobuf parser uses an intermediate representation to handle Protobuf.

A protobuf message on the wire is a sequence of records.

\vspace{1em}
\begin{definition}[Protobuf Record]
  A Protobuf wire record is a tuple with the field number $k$, the wire type $w$
  and the payload. The wire type will control how long the payload is.
\end{definition}

\vspace{1em}
\begin{definition}[Protobuf Record Wire Type]
  There are four Protobuf wire types represented in a record:
  \begin{itemize}
  \item \texttt{VARINT}, whose payload is an unsigned number up to $2^{64} - 1$.
  \item \texttt{I64}, whose payload is exactly 8 bytes.
  \item \texttt{I32}, whose payload is exactly 4 bytes.
  \item $\mathtt{LEN}\ \ell$, whose payload is exactly $\ell$ bytes.
  \end{itemize}
\end{definition}

% Local Variables:
% citar-bibliography: ("../pollux.bib")
% TeX-master: "../pollux.tex"
% TeX-command-default: "Make"
% jinx-local-words: "cardinality protobuf"
% End:
